\subsection{An exact class-number criterion at every finite abelian norm level}

\begin{claim}{B.AB1}{Proved}{Finite abelian norm obstructions, including nonconstant deck groups}
Let $D\to C$ be a pointed geometrically connected finite
\emph{\'etale} cover over $k=\mathbf F_q$, with abelian geometric
Galois group $A$ of order prime to $p=\operatorname{char}k$.
The chosen point of $D$ lies rationally above $P$.
Let $F$ act on $A$ by the pointed arithmetic descent, and put
\[
 S_D(k)=\{a\in J_D(k)[p^\infty]:
    \operatorname{Nm}_{D/D_H}(a)
    \text{ has trivial stabilizer in }A/H
                   \text{ for every }H\le A\},                 \tag{B.AB.1}
\]
where $D_H=D_{\overline k}/H$, with its pointed descent whenever
the subgroup is stable under a power of $F$.

For each subgroup $H$ for which $A/H$ is a nontrivial cyclic
prime-power group $\mathbf Z/\ell^a$, let $m_H$ be its orbit
length under $F$, and let $H^+/H$ be the unique subgroup of
order $\ell$ in $A/H$. Both quotient curves have their pointed
models over $k_{m_H}=\mathbf F_{q^{m_H}}$. Then
\[
 S_D(k)\ne\varnothing
 \quad\Longleftrightarrow\quad
 p\mid
 \frac{\#J_{D_H}(k_{m_H})}{\#J_{D_{H^+}}(k_{m_H})}
 \quad\text{for every such }H.                                \tag{B.AB.2}
\]
Thus an empty set at a finite abelian normal level is equivalent
to a relative class-number unit in one cyclic prime-power tower,
at that quotient's first pointed field of definition. The step
$D_H\to D_{H^+}$ has prime degree $\ell$; its lower curve need
not be $C$.
\end{claim}
\begin{proof}
We use the geometric $p$-primary torsion group
\[
                   T=J_D(\overline k)[p^\infty].
\]
For a subgroup $H\le A$, the element
\[
                   e_H=\frac1{|H|}\sum_{h\in H}h
                         \in\mathbf Z_p[A]
\]
is an idempotent. Pullback from $D_H$ is injective on
$p$-primary torsion, and pullback followed by norm is
$\sum_{h\in H}h=|H|e_H$. Its image is $e_HT$; this follows
also by applying norm and dividing by $|H|$. Consequently
the stabilizer in $A$ of $e_Ha$ is the inverse image of the
stabilizer of the norm coordinate in $A/H$. Freeness at this
quotient means exactly
\[
                         \operatorname{Stab}_A(e_Ha)=H.         \tag{B.AB.3}
\]

Here is the needed character decomposition, including its
integrality. Over a finite unramified extension of $\mathbf Z_p$
containing all roots of unity of exponent $\exp(A)$, the
idempotents of the characters $\chi$ diagonalize $A$.
For each subgroup $K$ with $A/K$ cyclic, sum these idempotents
over all characters whose kernel is exactly $K$, and denote
the result by $\varepsilon_K$. That sum is invariant under
the unramified Galois group and hence belongs to
$\mathbf Z_p[A]$. The $\varepsilon_K$ are pairwise orthogonal
and sum to one. Moreover
\[
       e_H=\sum_{K\supset H}\varepsilon_K.                     \tag{B.AB.4}
\]
Every nonzero element of $\varepsilon_KT$ has stabilizer
exactly $K$. To see this after the faithful unramified scalar
extension, if $g\notin K$, every corresponding value
$\chi(g)-1$ is a unit: a nontrivial root of unity of order
prime to $p$ remains nontrivial in the residue field. Thus
$g-1$ is invertible on that entire summand. This reasoning
holds for arbitrary finite $p$-power torsion, not only for
torsion killed by $p$.

If $A/H$ is cyclic of order $\ell^a$, its primitive packet is
particularly simple:
\[
                         \varepsilon_H=e_H-e_{H^+}.            \tag{B.AB.5}
\]
Every strict overgroup of $H$ contains $H^+$. Hence (B.AB.3)
forces $\varepsilon_Ha\ne0$: otherwise the whole norm
coordinate would be fixed by $H^+/H$. Conversely, suppose
$a$ has nonzero primitive component at \emph{every} subgroup
with cyclic prime-power quotient. For an arbitrary $K\le A$
and $g\notin K$, a prime-primary character of the finite
abelian group $A/K$ detects $g$. Its kernel lifts to a subgroup
$H\supset K$ with $A/H$ cyclic of prime-power order and
$g\notin H$. The nonzero component $\varepsilon_Ha$ survives
in $e_Ka$, and $g$ cannot fix it. Thus (B.AB.3) holds for
every $K$. We have proved the exact equivalence
\[
 a\text{ satisfies all norm-freeness conditions}
 \quad\Longleftrightarrow\quad
 \varepsilon_Ha\ne0\text{ for every cyclic prime-power }A/H.
                                                                  \tag{B.AB.6}
\]

We next impose Frobenius, without assuming that the deck group
is constant. Conjugation carries $\varepsilon_H$ to
$\varepsilon_{F(H)}$. If $a\in T^F=J_D(k)[p^\infty]$,
then $\varepsilon_Ha$ is fixed by $F^{m_H}$. Necessity in
(B.AB.6) therefore requires
\[
                         (\varepsilon_HT)^{F^{m_H}}\ne0.        \tag{B.AB.7}
\]
Conversely assume (B.AB.7) for every indicated $H$.
Choose one representative in each Frobenius orbit of these
subgroups and a nonzero $b_H$ in the displayed fixed group.
Then
\[
                  a_H=\sum_{i=0}^{m_H-1}F^i b_H
\]
is fixed by $F$. Its components in the distinct packets
$\varepsilon_{F^i(H)}T$ are all nonzero. Summing $a_H$ over
the finitely many subgroup orbits gives a rational torsion
class satisfying (B.AB.6). No division by $m_H$ occurs;
the proof is valid even when $p\mid m_H$.

Finally both $H$ and $H^+$ are stable under $F^{m_H}$, and
their quotient maps and rational lifts of $P$ descend to
$k_{m_H}$. Put
\[
 B_H=\bigl(\ker(J_{D_H}\longrightarrow J_{D_{H^+}})\bigr)^0.
\]
Pullback identifies $B_H[p^\infty]$ with the primitive packet
$\varepsilon_HT$, compatibly with $F^{m_H}$. Norm and
pullback split the $p$-primary groups because their degree
$\ell$ is prime to $p$. The isogeny
$J_{D_{H^+}}\times B_H\to J_{D_H}$ shows
\[
 \#B_H(k_{m_H})
       =\frac{\#J_{D_H}(k_{m_H})}{\#J_{D_{H^+}}(k_{m_H})}.
\]
Thus (B.AB.7) is exactly divisibility by $p$ of the indicated
integer. This proves (B.AB.2).
\end{proof}

\begin{claim}{B.AB2}{Proved}{The elementary-prime case and the stronger higher-power test}
If the group $A$ in \statementref{B.AB1} is constant, then all
$m_H=1$. If in addition $A=(\mathbf Z/\ell)^d$, then
$S_D(k)$ is nonempty exactly when every degree-$\ell$ quotient
$D_H\to C$ has relative class number divisible by $p$.
Thus at a constant elementary-abelian level, an obstruction
is exactly a prime-degree class-number witness among those
quotients.

For higher-exponent abelian groups, the exact criterion also
tests the primitive parts at cyclic prime-power quotients.
These tests are strictly stronger as finite module conditions
than testing only prime-degree quotients. This strictness does
not assert that a particular curve has no prime-degree witness
over all constant extensions.
\end{claim}
\begin{proof}
For elementary $A$, every nontrivial cyclic quotient has degree
$\ell$, and $H^+=A$. The relative factor in (B.AB.2) is
therefore exactly that of $D_H/C$. This proves the first
assertion and identifies the relationship with B1 at the
displayed finite field and level.

Here is a complete finite-module control for the higher-power
assertion. Take characteristic two and
\[
        A=\mathbf Z/9\times\mathbf Z/3,
                     \qquad M=\mathbf F_4^4.
\]
Let $\omega\in\mathbf F_4$ have order three. An element
$(a,b)$ acts diagonally on the four summands with scalars
\[
          \omega^a,\qquad\omega^b,\qquad
          \omega^{a+b},\qquad\omega^{a+2b}.                    \tag{B.AB.8}
\]
These are all four character kernels of prime-degree
quotients. A vector with all four coordinates nonzero has
free norm orbit at every such quotient. There are $3^4=81$
such vectors. Nevertheless the subgroup
\[
                   \{(0,0),(3,0),(6,0)\}
\]
fixes every element of $M$, so no vector has free orbit at
the full level. The three missing primitive packets have
cyclic quotient of order nine.

For comparison, choose a primitive ninth root
$\zeta\in\mathbf F_{64}$ and add three summands
$\mathbf F_{64}$ with characters
\[
                \zeta^{a+3jb},\qquad j=0,1,2.
\]
Now all seven cyclic prime-power kernels are present.
The vector with every coordinate one has free norm orbit
at all ten subgroups of $A$. In fact the exact number of
all-quotient free vectors in this module is
\[
                       3^4\cdot63^3=20\,253\,807.
\]
The packet criterion proves the formula; the accompanying
program checks every subgroup and every norm for the displayed
vector. It exhausts the $256$ vectors of the first module and
independently verifies its count $81$ and its absence of an
all-quotient free vector.

There is also a check of the nonconstant Frobenius step when
the orbit length is divisible by the characteristic. Set
$A=(\mathbf Z/3)^2$, retain (B.AB.8), and define
\[
 F(a,b)=(b,a),\qquad
             F(u,v,w,z)=(v,u,w,z^2).
\]
These actions satisfy $F(gm)=F(g)F(m)$. The fixed module
has $32$ vectors, of which exactly nine have free norm orbit
at every subgroup. The first two character kernels form an
orbit of length two, and the other two are stable. The sum
\[
                  (1,0,0,0)+F(1,0,0,0)=(1,1,0,0)
\]
is fixed and is nonzero in both moving packets, despite
working in characteristic two. Thus the Frobenius-orbit
argument in \statementref{B.AB1} does not hide an averaging
division by an inadmissible integer.

All finite-field operations, subgroup enumeration, norms,
stabilizers and counts are independently implemented in
\path{code/criterion_abelian_norm.py}; its complete results
are in the adjacent JSON certificate. These are finite
representation examples, with no claimed Jacobian realization.
The universal genus-two characteristic-two assertion in B2
still requires an empty level for every pointed curve. The
present theorem exactly reduces every \emph{abelian} candidate
level to the displayed relative class-number tests; it does
not settle the existence of such a level, or rule out a
stronger obstruction at a nonabelian level.
\end{proof}
